\section{Dependency graph and the role of dynamics}
\label{sec:dependency}

\subsection{The full logical graph}

The route is intentionally drawn with conjectural nodes visible.

\medskip
\begin{center}
\fbox{\begin{minipage}{0.92\linewidth}
\centering
\(\Etag\) exact fixed-shift local-model identity
\(\eqref{eq:hard-packet-anchor}\)\\[0.4em]
\(\Uparrow\)\\[-0.2em]
\(\Ctag\) original hard packet \(B_{h_0,\delta}(X)=o(X)\)\\[0.4em]
\(\Uparrow\)\\[-0.2em]
\(\Htag\) H8 literal reconnection
\(\quad+\quad\)
\(\Htag\) H9 strict endpoint\\[0.4em]
\(\Uparrow\)\\[-0.2em]
\(\Htag\) H6 full-block return
\(\quad+\quad\)
\(\Htag\) H7 fixed-\(h_0\) localization\\[0.4em]
\(\Uparrow\)\\[-0.2em]
\(\Htag\) H1 canonical physical synthesis\\[0.5em]
\(\Uparrow\)\\[-0.2em]
\begin{tabular}{c@{\qquad\qquad}c}
positive low-frequency resonance & generic/high sectors\\
\(\Etag\) two exceptional returns & \(\Htag\) H3 signed affine arithmetic\\
\(\Etag\) three-ledger inequality & \(\Htag\) H4 high/ultra return\\
\(\Htag\) H2 \(W/\mathfrak P/\mathfrak X\) &
literal outer coefficients
\end{tabular}\\[0.5em]
\(\Downarrow\)\\[-0.2em]
\(\Htag\) H5 determinant-energy/zero-mode compatibility
\end{minipage}}
\end{center}
\medskip

The graph has two important forks.

\begin{enumerate}[leftmargin=2em]
\item H2 is a \emph{positive} sufficient route.  If a literal
  principal or cross-map ledger is too large, the route may switch to
  a signed filter before taking absolute values.  Failure of H2 does
  not disprove the scalar correlation.
\item H3 may be approached through direct native-\(h_0\) determinant
  dispersion, a coherent prime-square theorem, or a signed affine
  large sieve.  A result on an averaged shift or random coefficient
  is a different branch, not a proof of H3.
\end{enumerate}

\subsection{Where the parity difficulty lives}

The program has already classified the local squarefree support in
the primitive affine sector.  After that classification the remaining
sign is an actual growing Liouville/M\"obius parity correlation
\cite{WangTPC90}.  Accordingly, the architecture does not claim that
the parity barrier has been partly broken.  It claims something more
modest and testable:
\begin{quote}
the literal carrier and normalization at which parity-sensitive
cancellation must occur have been isolated more sharply.
\end{quote}
H3 is the node where an actual arithmetic advance must enter.  A
support theorem, positive density, spectral decomposition, or finite
certificate cannot substitute for it.

\subsection{What the dynamical language contributes}

The dynamical and spectral language used throughout the broader
program contributes:
\begin{itemize}[leftmargin=2em]
\item canonical symbolic, orbit, conductor, and affine coordinates;
\item transfer-like decompositions and exact finite resonance modes;
\item Gram, projection, energy, holonomy, and width-sensitive norm
  identities;
\item branch classification, provenance, and loss ledgers; and
\item sharp countermodels that identify where a proposed transfer
  loses the physical zero mode.
\end{itemize}
These are genuine structural uses.  In particular, the multiplicative
character spectrum behind \eqref{eq:three-ledger} replaces a crude
worst-eigenvalue estimate by the exact filter factor \(g_q(H)\)
\cite{WangTPC99,WangTPC101}.

\subsection{What it does not contribute automatically}

No dynamical or finite spectral identity in this program automatically
supplies:
\begin{itemize}[leftmargin=2em]
\item independence of primes;
\item M\"obius randomness on the actual fixed-\(h_0\) carrier;
\item a fixed-shift arithmetic cancellation theorem;
\item a breach of sieve parity; or
\item an identification of Riemann zeros with eigenvalues of the TPC
  operators.
\end{itemize}
The safe summary is:
\begin{quote}
\emph{Dynamics organizes the arithmetic carrier and exposes where
cancellation must occur; it does not itself supply that
cancellation.}
\end{quote}
No statement in TPC-MVP1 concerns the Riemann hypothesis,
Hilbert--P\'olya, or a zeta-spectrum realization.

\subsection{Two exponent geometries}

The current packet scales
\[
Q_X=X^{267/400+o(1)},
\qquad
J_X=X^{133/400+o(1)}
\]
satisfy \(Q_XJ_X=X^{1+o(1)}\).  On a resolved long fiber
\(N\ge J_X\), the established width-sensitive resonance norm replaces
the crude exponent \(267/800\) by \(67/400\), a gain
\[
\frac{267}{800}-\frac{67}{400}
=\frac{133}{800}.
\]
This is a real native spectral improvement on those fibers.  It is
not a global packet estimate because short fibers, principal mass,
cross-map excess, and physical synthesis still have to be returned.

Independently, the physical synthesis has only the strict allowance
\[
\Lambda_{\mathrm{phys}}<\frac1{400}.
\]
The large local numerical gain above may not be spent twice or
transferred to unrelated blocks without a theorem.
